\chapter{Statistiques d'ordre}\label{ch:iv1:order-statistics}

\begin{ivtrains}
Travailler avec les valeurs triées d'un échantillon~: maximum et médiane, écarts entre valeurs successives, plus grand numéro de série observé, deuxième meilleur cours acheteur, records. Les rangs sont le fil conducteur, car la configuration des rangs d'un échantillon ne dépend pas de sa loi.
\end{ivtrains}

\begin{question}[$\star$ \iqroles{QR, QT, Trader, Analyste} \iqfirm{toutes} \iqlevel{junior}]\label{iq:iv1:order-statistics:1}
Vous tirez $n$ nombres indépendants uniformes sur $[0, 1]$. Quelles sont les espérances du plus grand, du plus petit et de l'étendue~? Donnez la probabilité que le plus grand soit inférieur à $0.9$ lorsque $n = 10$.
\end{question}

\begin{question}[$\star$ \iqroles{QR, QT, DS, Analyste} \iqfirm{toutes} \iqlevel{junior}]\label{iq:iv1:order-statistics:2}
Vous tirez trois nombres indépendants uniformes sur $[0, 1]$. Quelle est la densité de leur médiane~? Quelles sont son espérance et sa variance~? Lequel estime mieux le centre de la distribution, la médiane des trois nombres ou leur moyenne~?
\end{question}

\begin{question}[$\star$ \iqroles{QR, QT, Trader, DS} \iqfirm{toutes} \iqlevel{junior}]\label{iq:iv1:order-statistics:3}
Le desk A envoie 10 ordres et le desk B en envoie 20, et les 30 latences sont des tirages indépendants d'une même loi continue. Quelle est la probabilité que l'ordre le plus rapide appartienne au desk B~? Que les trois plus rapides appartiennent tous au desk B~? La réponse dépend-elle de la loi~?
\end{question}

\begin{question}[$\star\star$ \iqroles{QR, QT, DS, Analyste} \iqfirm{toutes} \iqlevel{junior}]\label{iq:iv1:order-statistics:4}
Un lot de machines porte les numéros de série $1, 2, \dots, N$, où $N$ est inconnu. Vous observez cinq machines différentes, tirées au hasard, et le plus grand numéro de série parmi elles est 60. Donnez une estimation sans biais de $N$, et démontrez-la. Pourquoi le double de la moyenne de l'échantillon moins un est-il un moins bon estimateur, alors qu'il est lui aussi sans biais~?
\end{question}

\begin{question}[$\star\star$ \iqroles{QR, QT, DS} \iqfirm{toutes} \iqlevel{junior}]\label{iq:iv1:order-statistics:5}
Vous disposez de $n$ tirages indépendants d'une loi uniforme sur $[0, \theta]$. Comparez la variance de $2\bar{X}$ à celle de $\frac{n + 1}{n}\max X_i$ pour $n = 10$. Les deux estimateurs sont sans biais. Quel multiple du maximum a le plus petit erreur quadratique moyenne, et quelle est sa valeur~?
\end{question}

\begin{question}[$\star\star$ \iqroles{QR, QT, Trader, Analyste} \iqfirm{teneur de marché, hedge fund} \iqlevel{junior}]\label{iq:iv1:order-statistics:6}
Cinq cours acheteurs sont indépendants et uniformes sur $[0, 1]$. Quelle est l'espérance du deuxième cours acheteur le plus élevé~? Si l'on vous indique que le cours acheteur le plus élevé vaut $0.8$, quelle est alors cette espérance~? Par ailleurs, vous lancez quatre dés équilibrés~: quelle est l'espérance du deuxième plus grand chiffre obtenu~?
\end{question}

\begin{question}[$\star\star$ \iqroles{QR, QT, DS, Risque} \iqfirm{toutes} \iqlevel{junior}]\label{iq:iv1:order-statistics:7}
Pour un grand échantillon issu d'une loi normale, de combien la médiane de l'échantillon est-elle plus variable que sa moyenne~? Quelle est la même comparaison pour la loi de Laplace, de densité $\frac12 e^{-|x|}$~? Qu'est-ce que cela suggère pour les rendements à queues épaisses~?
\end{question}

\begin{question}[$\star\star$ \iqroles{QR, QT, Risque, QD} \iqfirm{toutes} \iqlevel{junior}]\label{iq:iv1:order-statistics:8}
Cinq serveurs tombent en panne après des durées exponentielles indépendantes, de moyenne un an. En moyenne, au bout de combien de temps le premier tombe-t-il en panne~? Le troisième~? Le dernier~? Le temps de la première panne et le temps écoulé entre la première et la dernière panne sont-ils indépendants~?
\end{question}

\begin{question}[$\star\star$ \iqroles{QR, QT, DS} \iqfirm{toutes} \iqlevel{junior}]\label{iq:iv1:order-statistics:9}
Vous tirez des nombres indépendants d'une loi continue jusqu'à ce que l'un d'eux dépasse votre premier tirage. Soit $N$ le nombre de tirages supplémentaires nécessaires. Quelle est $P(N > n)$~? Quel est le plus petit $n$ tel que $P(N \le n) \ge \frac12$, et quelle est l'espérance de $N$~? Les événements «~le tirage 5 est un record~» et «~le tirage 10 est un record~» sont-ils indépendants~?
\end{question}

\begin{question}[$\star\star$ \iqroles{QR, QT, DS, Risque} \iqfirm{toutes} \iqlevel{junior}]\label{iq:iv1:order-statistics:10}
Vous disposez de 10 tirages indépendants d'une loi continue inconnue. Quelle est la probabilité que sa vraie médiane se situe entre le plus petit et le plus grand tirage~? Entre le troisième plus petit et le troisième plus grand~? Pourquoi peut-on l'affirmer sans rien savoir de la loi~?
\end{question}

\begin{question}[$\star\star\star$ \iqroles{QR, QT, DS} \iqfirm{hedge fund, teneur de marché} \iqlevel{junior, confirmé}]\label{iq:iv1:order-statistics:11}
Un chercheur effectue le backtest de 1000 stratégies indépendantes, toutes de ratio de Sharpe vrai égal à 0. Chaque ratio de Sharpe annuel mesuré suit une loi normale centrée en 0, d'écart type $1/3$, ce qui correspond à 9 années de données. Quelle est, à peu près, la meilleure valeur mesurée du ratio de Sharpe, sans recourir à une table~? Comment l'espérance de la meilleure valeur croît-elle avec le nombre de stratégies testées~?
\end{question}

\begin{question}[$\star\star\star$ \iqroles{QR, QT, Risque, DS} \iqfirm{toutes} \iqlevel{junior, confirmé}]\label{iq:iv1:order-statistics:12}
Vous tirez $n$ nombres indépendants uniformes sur $[0, 1]$. Déterminez la loi de l'étendue. Quelle doit être la valeur de $n$ pour que la probabilité que l'étendue couvre au moins 90~\% de l'intervalle soit d'au moins 95~\%~? Que peut-on alors affirmer, avec ce même $n$, pour une loi continue quelconque~?
\end{question}

\begin{question}[$\star\star\star$ \iqroles{QR, QT, DS} \iqfirm{hedge fund} \iqlevel{junior, confirmé}]\label{iq:iv1:order-statistics:13}
Un échantillon de taille 11 provient d'une loi uniforme sur $[\theta, \theta + 1]$, où $\theta$ est inconnu. Comparez trois estimateurs du centre~: la moyenne de l'échantillon, la médiane de l'échantillon et le milieu de l'étendue, c'est-à-dire la moyenne du plus petit et du plus grand. Donnez les trois variances. Pourquoi le milieu de l'étendue l'emporte-t-il ici, et à quelle vitesse sa variance décroît-elle lorsque l'échantillon grandit~?
\end{question}

\begin{question}[$\star\star\star$ \iqroles{QR, QT, LLE, QD} \iqfirm{teneur de marché, crypto} \iqlevel{junior, confirmé}]\label{iq:iv1:order-statistics:14}
Dix ordres arrivent à des instants uniformes et indépendants au cours d'une seconde. Quelle est la probabilité qu'aucun couple d'ordres ne soit distant de moins de 10 millisecondes~? Quelle est l'espérance du plus petit écart entre deux ordres voisins~?
\end{question}

\begin{question}[$\star\star\star\star$ \iqroles{QR, QT} \iqfirm{hedge fund, teneur de marché} \iqlevel{confirmé}]\label{iq:iv1:order-statistics:15}
Calculez, sous forme fermée, l'espérance du plus grand de trois variables normales centrées réduites indépendantes, puis celle du plus grand de quatre. Votre méthode ne doit pas nécessiter d'intégration numérique. Jusqu'où va la même méthode, et qu'est-ce qui l'arrête~?
\end{question}
